(1) 单色激光真空波长 $\lambda$ 进入理想迈克耳孙干涉仪。两臂回程光强为 $I_1,I_2$，复相干度模为 $\mathcal V$，求输出光强随两臂相差 $\phi$ 的表达式。
(2) 将其中一面反射镜沿光轴移动 $\Delta x$，求由此产生的相位变化及经过的完整条纹数。
(3) 在一臂单程放入折射率 $n>1$、厚度 $d$ 的透明平板，另一臂仍为空气，忽略色散和反射损失，求附加光程差与相位变化。
(4) 给出测镜位移的相位灵敏度，讨论不等臂光强及部分相干对条纹可见度的影响。
